\documentclass[12pt]{article}
\usepackage[utf8]{inputenc}
\usepackage[T1]{fontenc}
\usepackage{hyperref}
\title{Perai Language Resources for K-12 Curriculum Development}
\author{Compiled by Claude Sonnet 5.5 (claude-sonnet-5-5) for Andr\'as Kornai}
\date{}
\begin{document}
\maketitle

\section{Introduction}
Perai (ISO 639-3 \texttt{wet}; Glottocode pera1257) is described in the task as an endangered Austronesian language of Indonesia with about 280 speakers, used as a first language by older adults only. My searches found no source about Perai. The results concerned other languages, such as Alorese, and nothing connected them to Perai.

\section{Primary Educational Resources}
None found.

\section{Linguistic and Reference Materials}
None found. I could not confirm any grammar, dictionary or word list for Perai.

\section{Religious and Community Resources}
None found.

\section{Digital and Audio Resources}
None found.

\section{Curriculum Development Context}
With about 280 speakers, all older adults, the language needs urgent documentation before any school curriculum is feasible. Fieldwork and recordings with the remaining speakers would come first. I found no existing programme.
\end{document}